\paragraph{\texorpdfstring{$S_4$}{S4} 闭包下的 Chevalley--Weil 重数刚性、Jacobian 群代数分解与局部算术证书}
沿用结论 \ref{con:xi-terminal-zm-s4-galois-closure-genus-tower-holomorphic-differentials} 的设定，
记 \(\widetilde{X}\to \PP^1_y\) 为四次谱覆盖 \(\Pi(\lambda,y)=0\) 的 \(S_4\)-伽罗瓦闭包曲线，且 \(g(\widetilde{X})=16\)。
并记
\[
F(\lambda,y):=\Pi(\lambda,y)
=\lambda^{4}-\lambda^{3}-(2y+1)\lambda^{2}+\lambda+y(y+1).
\]
记
\[
V:=H^0(\widetilde{X},\Omega^1_{\widetilde{X}}),
\qquad
\varepsilon:=\mathrm{sgn},\ \rho_2:=V_2,\ \rho_3:=V_3,\ \rho_3':=V_3'.
\]

\begin{conclusion}[\texorpdfstring{$S_4$}{S4} 闭包覆盖的 Hurwitz 子群谱：全部中间商曲线亏格与惯性循环型]\label{con:xi-terminal-zm-s4-hurwitz-subgroup-spectrum-cycletypes}
设 \(\tau=(12)\)、\(\rho=(1234)\in S_4\)。
对任意子群 \(H\le S_4\)，记
\[
d:=[S_4:H],\qquad
c_\tau(H):=\#\text{\(\tau\) 在 }S_4/H\text{ 上的循环个数},\qquad
c_\rho(H):=\#\text{\(\rho\) 在 }S_4/H\text{ 上的循环个数}.
\]
则有统一亏格公式
\[
g(\widetilde{X}/H)
=1-d+\frac12\Bigl(5\bigl(d-c_\tau(H)\bigr)+\bigl(d-c_\rho(H)\bigr)\Bigr).
\]
并且在各共轭类代表子群上，可得完整枚举
\[
\begin{array}{c|c|c|c|c}
H & d & \tau\text{ 在 }S_4/H\text{ 上的循环型} & \rho\text{ 在 }S_4/H\text{ 上的循环型} & g(\widetilde{X}/H)\\
\hline
S_4 & 1 & 1 & 1 & 0\\
A_4 & 2 & 2 & 2 & 2\\
S_3 & 4 & 2\cdot 1^2 & 4 & 1\\
D_8 & 3 & 2\cdot 1 & 2\cdot 1 & 1\\
V_4^{\mathrm{norm}} & 6 & 2^3 & 2^3 & 4\\
V_4^{\mathrm{pair}}=\langle(12),(34)\rangle & 6 & 2^2\cdot 1^2 & 4\cdot 2 & 2\\
C_4 & 6 & 2^3 & 4\cdot 1^2 & 4\\
C_3 & 8 & 2^4 & 4^2 & 6\\
C_2^{\mathrm{tr}} & 12 & 2^5\cdot 1^2 & 4^3 & 6\\
C_2^{\mathrm{dt}} & 12 & 2^6 & 4^2\cdot 2^2 & 8\\
\{1\} & 24 & 2^{12} & 4^6 & 16
\end{array}
\]
其中 \(D_8\) 表示阶 \(8\) 二面体子群，\(V_4^{\mathrm{norm}}\) 为正规 Klein 四元，\(V_4^{\mathrm{pair}}\) 为两两换位 Klein 四元。
\end{conclusion}
\begin{proof}[证明要点]
\(\widetilde{X}\to\PP^1_y\) 的惯性共轭类为五个换位与一个四循环。
对商覆盖 \(\widetilde{X}/H\to\PP^1_y\)，由置换 \(g\in\{\tau,\rho\}\) 在 \(S_4/H\) 上的循环分解可读出总分歧量 \(d-c_g(H)\)；
代入 Riemann--Hurwitz 公式即得统一亏格表达式。
表中循环型由 \(S_4\) 在左陪集上的显式置换表示计算给出。证毕。
\end{proof}

\begin{conclusion}[指数 \texorpdfstring{$6$}{6} 配对稳定子商曲线的显式模型]\label{con:xi-terminal-zm-s4-v4pair-explicit-model}
在代数闭包中记 \(F(\lambda,y)=0\) 的四根为 \(\lambda_1,\lambda_2,\lambda_3,\lambda_4\)，并取
\[
V_4^{\mathrm{pair}}=\langle(12),(34)\rangle\le S_4.
\]
定义
\[
v:=\lambda_1\lambda_2,\quad w:=\lambda_3\lambda_4,\quad z:=v+w,\quad s:=v-w.
\]
则固定域可写为
\[
\QQ(\widetilde{X}/V_4^{\mathrm{pair}})=\QQ(y,z,s),
\]
并满足完全显式关系
\[
\mathscr R(z,y)=0,\qquad s^2=z^2-4y(y+1),
\]
其中
\[
\mathscr R(z,y)=z^3+(1+2y)z^2-(1+4y+4y^2)z-(1+5y+13y^2+8y^3).
\]
记 \(E_{\mathrm{res}}\) 为 \(\mathscr R(z,y)=0\) 的光滑射影模型，\(C_{\mathrm{pair}}\) 为上述二元方程组的光滑射影模型，则
\[
\QQ(C_{\mathrm{pair}})=\QQ(\widetilde{X}/V_4^{\mathrm{pair}}),\qquad
g(C_{\mathrm{pair}})=2,
\]
且自然投影 \(C_{\mathrm{pair}}\to E_{\mathrm{res}}\) 为次数 \(2\) 覆盖。
\end{conclusion}
\begin{proof}[证明要点]
\(z=v+w\) 为配对 \(\{12|34\}\) 的不变量，故满足三次预解式 \(\mathscr R(z,y)=0\)。
又由 \(vw=\lambda_1\lambda_2\lambda_3\lambda_4=y(y+1)\)，得
\[
s^2=(v-w)^2=(v+w)^2-4vw=z^2-4y(y+1).
\]
于是 \(\QQ(y,z,s)\subseteq \QQ(\widetilde{X}/V_4^{\mathrm{pair}})\)。
反向由 \(v=(z+s)/2,\ w=(z-s)/2\) 可恢复配对乘积信息，故函数域相等。
亏格与覆盖次数由结论 \ref{con:xi-terminal-zm-s4-hurwitz-subgroup-spectrum-cycletypes} 中 \(H=V_4^{\mathrm{pair}}\) 的条目给出。证毕。
\end{proof}

\begin{conclusion}[黄金比率分歧：\texorpdfstring{$C_{\mathrm{pair}}\to E_{\mathrm{res}}$}{Cpair to Eres} 的分歧除子由 \texorpdfstring{$y^2+y-1$}{y2+y-1} 唯一刻画]\label{con:xi-terminal-zm-s4-v4pair-golden-ratio-branch-divisor}
在结论 \ref{con:xi-terminal-zm-s4-v4pair-explicit-model} 的记号下，令
\[
f:=z^2-4y(y+1)\in \QQ(E_{\mathrm{res}})^\times.
\]
则有严格消去恒等式
\[
\operatorname{Res}_z\!\bigl(\mathscr R(z,y),\,z^2-4y(y+1)\bigr)=(y^2+y-1)^2.
\]
因此 \(f=0\) 与 \(\mathscr R=0\) 的公共点恰为
\[
P_\pm:\quad
(y,z)=\Bigl(-\frac12\pm\frac{\sqrt5}{2},\,-2\Bigr),
\]
并且二次覆盖 \(C_{\mathrm{pair}}\to E_{\mathrm{res}}\) 的分歧除子为 \(P_++P_-\)。
从而
\[
2g(C_{\mathrm{pair}})-2=\deg(P_++P_-)=2,\qquad g(C_{\mathrm{pair}})=2,
\]
且分歧点由唯一二次域 \(\QQ(\sqrt5)\) 控制。
\end{conclusion}
\begin{proof}[证明要点]
结果式计算给出仅当 \(y^2+y-1=0\) 时可能有公共点。
将该二次方程两根代回 \(\mathscr R(z,y)=0\) 与 \(z^2-4y(y+1)=0\) 可唯一得到 \(z=-2\)。
对二次覆盖 \(s^2=f\)，分歧由 \(f\) 的奇阶零点给出，故分歧除子正为 \(P_++P_-\)。
再由 Riemann--Hurwitz 公式即得亏格。证毕。
\end{proof}

\begin{conclusion}[Chevalley--Weil 型 \texorpdfstring{$S_4$}{S4} 表示分解的显式重数解]\label{con:xi-terminal-zm-s4-chevalley-weil-explicit-decomposition}
在上述记号下，\(V\) 作为 \(S_4\)-表示满足
\[
V\cong 2\varepsilon\oplus \rho_2\oplus \rho_3\oplus 3\rho_3'.
\]
等价地，\(\dim V=16\) 的不可约重数向量唯一为
\[
(m_{\varepsilon},m_{\rho_2},m_{\rho_3},m_{\rho_3'})=(2,1,1,3).
\]
\end{conclusion}
\begin{proof}[证明要点]
由结论 \ref{con:xi-terminal-zm-s4-galois-closure-genus-tower-holomorphic-differentials}，
\[
g(\widetilde{X}/S_4)=0,\quad g(\widetilde{X}/A_4)=2,\quad g(\widetilde{X}/V_4)=4,\quad g(\widetilde{X}/S_3)=1.
\]
对任意子群 \(K\le S_4\) 有 \(\dim V^K=g(\widetilde{X}/K)\)。
设 \(V\cong a\mathbf 1\oplus b\varepsilon\oplus c\rho_2\oplus d\rho_3\oplus e\rho_3'\)。
由 \(\dim V^{S_4}=0\) 得 \(a=0\)；再由角色平均公式给出的不变维数
\[
\dim(\varepsilon^{A_4})=1,\quad \dim(\rho_2^{V_4})=2,\quad \dim(\rho_3^{S_3})=1,
\]
以及其余对应项为零，得到
\[
b=2,\qquad b+2c=4,\qquad d=1,
\]
故 \(c=1\)。
最后由 \(\dim V=2b+2c+3d+3e=16\) 解得 \(e=3\)。
唯一性由线性方程组满秩给出。证毕。
\end{proof}

\begin{conclusion}[Jacobian 的群代数分解与中间 Jacobian 的显式同源分拆]\label{con:xi-terminal-zm-s4-jacobian-group-algebra-prym-e3-b3}
记
\[
J:=J(\widetilde{X}),\qquad
J_{\mathrm{LY}}:=J(\widetilde{X}/A_4),\qquad
E:=\widetilde{X}/S_3,\qquad
E_{\mathrm{res}}:=\widetilde{X}/D_8.
\]
则存在定义在 \(\QQ\) 上的三维阿贝尔簇 \(B\)，使得
\[
J\sim_{\QQ}J_{\mathrm{LY}}\times E_{\mathrm{res}}^{2}\times E^{3}\times B^{3}.
\]
等价地，也可写作
\[
J\sim_{\QQ}J(\widetilde{X}/A_4)\times \operatorname{Prym}\!\bigl((\widetilde{X}/V_4^{\mathrm{norm}})/(\widetilde{X}/A_4)\bigr)\times E^{3}\times B^{3}.
\]
并且以下中间 Jacobian 分拆成立：
\[
J(\widetilde{X}/V_4^{\mathrm{norm}})\sim_{\QQ}J_{\mathrm{LY}}\times E_{\mathrm{res}}^{2},
\]
\[
J(C_{\mathrm{pair}})\sim_{\QQ}E_{\mathrm{res}}\times E,\qquad
C_{\mathrm{pair}}:=\widetilde{X}/V_4^{\mathrm{pair}},
\]
\[
J(\widetilde{X}/C_4)\sim_{\QQ}E_{\mathrm{res}}\times B.
\]
在不出现额外内自同态的泛型情形下，可取 \(\mathrm{End}^0_{\overline{\QQ}}(B)=\QQ\)。
\end{conclusion}
\begin{proof}[证明要点]
由
\(
\QQ[S_4]\cong \QQ\oplus \QQ\oplus M_2(\QQ)\oplus M_3(\QQ)\oplus M_3(\QQ)
\)
与结论 \ref{con:xi-terminal-zm-s4-chevalley-weil-explicit-decomposition}，
可将 \(J\) 分解为 \(\varepsilon,\rho_2,\rho_3,\rho_3'\) 四个等型部分，其维数分别为 \(2,2,3,9\)。
\(\varepsilon\)-分量由 \(A_4\) 商识别为 \(J_{\mathrm{LY}}\)，\(\rho_3\)-分量由 \(S_3\) 商识别为 \(E^3\)。
由结论 \ref{con:xi-terminal-zm-s4-prym-c6-over-c-d8-elliptic-square}，\(\rho_2\)-分量同源于 \(E_{\mathrm{res}}^2\)（亦即对应 Prym 二维因子）。
\(\rho_3'\)-分量因重数为 \(3\) 可写作 \(B^3\)。
\par
对中间子群 \(V_4^{\mathrm{norm}},V_4^{\mathrm{pair}},C_4\)，由
\(
H^0(\widetilde{X}/H,\Omega^1)\cong H^0(\widetilde{X},\Omega^1)^H
\)
以及结论 \ref{con:xi-terminal-zm-s4-hurwitz-subgroup-spectrum-cycletypes} 的亏格数据，
再配合角色平均计算各不可约表示的 \(H\)-不变维数，即得三条同源分拆。证毕。
\end{proof}

\begin{conclusion}[\texorpdfstring{$D_8$}{D8} 商椭圆曲线与二维 Prym 的椭圆平方化]\label{con:xi-terminal-zm-s4-prym-c6-over-c-d8-elliptic-square}
设 \(D_8\le S_4\) 为任一 \(4\)-循环的正规化子（阶 \(8\) 二面体子群），并定义
\[
E_8:=\widetilde{X}/D_8.
\]
并记
\[
C:=\widetilde{X}/A_4,\qquad C_6:=\widetilde{X}/V_4^{\mathrm{norm}}.
\]
则 \(g(E_8)=1\)，并且
\[
\operatorname{Prym}(C_6/C)\sim_{\QQ}E_8^2.
\]
从而
\[
\mathrm{End}^0_{\QQ}\!\bigl(\operatorname{Prym}(C_6/C)\bigr)\supset M_2(\QQ).
\]
在无分歧循环三次覆盖情形下，\(\operatorname{Prym}(C_6/C)\) 的极化型为 \((1,3)\)。
\end{conclusion}
\begin{proof}[证明要点]
由结论 \ref{con:xi-terminal-zm-s4-galois-closure-genus-tower-holomorphic-differentials}，\(g(\widetilde{X}/D_8)=1\)。
由结论 \ref{con:xi-terminal-zm-s4-jacobian-a2-isog-r2}，\(\rho_2\)-等型分量满足 \(A_2\sim R^2\)，其中 \(R=\widetilde{X}/D_8\)。
另一方面，定理 \ref{thm:xi-terminal-zm-s3-closure-jacobian-splitting} 识别 \(A_2\sim \operatorname{Prym}(C_6/C)\)。
合并即得 \(\operatorname{Prym}(C_6/C)\sim_{\QQ}E_8^2\)。
极化型由定理 \ref{thm:xi-terminal-zm-prym-polarization-a2-rigidity} 给出。证毕。
\end{proof}

\begin{conclusion}[好素数处的局部 \texorpdfstring{$L$}{L} 因子强分解与模 \texorpdfstring{$3$}{3} 点数同余证书]\label{con:xi-terminal-zm-s4-local-factorization-mod3-congruence-certificate}
设 \(p\) 为 \(\widetilde{X}\) 的好约化素数，并令 \(q=p^m\)（\(m\ge 1\)）。
沿用结论 \ref{con:xi-terminal-zm-s4-jacobian-group-algebra-prym-e3-b3} 与
\ref{con:xi-terminal-zm-s4-prym-c6-over-c-d8-elliptic-square} 的记号，并记
\[
C:=\widetilde{X}/A_4.
\]
则局部因子满足
\[
P_p(\widetilde{X},T)
=P_p(C,T)\,P_p(E,T)^3\,P_p(E_8,T)^2\,P_p(B,T)^3.
\]
进一步，对有限域点数记号 \(N_Y(q):=\#Y(\FF_q)\)，有同余
\[
N_{\widetilde{X}}(q)\equiv N_C(q)+2N_{E_8}(q)-2(q+1)\pmod 3.
\]
\end{conclusion}
\begin{proof}[证明要点]
\(\QQ\)-同源分解在好约化处与 \(\ell\)-进 \(H^1\) 分解相容，故 Frobenius 特征多项式按因子相乘，得到 \(P_p\) 强分解。
又
\[
N_Y(q)=q+1-\mathrm{Tr}\!\left(\mathrm{Frob}_q\mid H^1_{\acute et}(Y_{\overline{\FF}_q},\QQ_\ell)\right).
\]
由分解式可得
\[
\mathrm{Tr}_q(\widetilde{X})
\equiv \mathrm{Tr}_q(C)+2\mathrm{Tr}_q(E_8)\pmod 3,
\]
因为 \(3\mathrm{Tr}_q(E)\) 与 \(3\mathrm{Tr}_q(B)\) 在模 \(3\) 下消失。
代回点数公式即得所示同余。证毕。
\end{proof}

\paragraph{公开检索说明}
对上述结构中的关键关键词组合（包括 \(S_4\)-闭包、表示重数向量、子群商属谱与 Prym 平方化识别）进行公开检索时，
未检得与本文同一覆盖对象逐项一致的组合性陈述。
该观察仅构成方法学层面的检索证据，不替代完备文献综述下的逻辑唯一性证明。
